\documentclass[a4paper,12pt]{extarticle}

\usepackage{cmap}
\usepackage{geometry}
\usepackage[T2A]{fontenc}
\usepackage[utf8]{inputenc}
\usepackage[english,russian]{babel}

\usepackage{amsmath}
\usepackage{amsthm}
\usepackage{amssymb}
\usepackage{mathtools}

\usepackage{setspace}
\usepackage{graphicx}
\usepackage{booktabs}
\usepackage{multirow}
\usepackage{tabularx}
\usepackage[flushleft]{threeparttable}
\usepackage{tablefootnote}
\usepackage{subcaption}
\usepackage{placeins}
\usepackage{float}
\usepackage{xurl}
\usepackage{indentfirst}
\graphicspath{{graphics/}}

\usepackage{chngcntr}
\counterwithin{table}{section}
\counterwithin{figure}{section}

\usepackage[numbers]{natbib}
\usepackage[colorlinks,citecolor=blue,linkcolor=blue,bookmarks=false,hypertexnames=true,urlcolor=blue]{hyperref}

\geometry{left=2.5cm}
\geometry{right=1.0cm}
\geometry{top=2.0cm}
\geometry{bottom=2.0cm}

\setlength{\parindent}{1.25cm}
\renewcommand{\baselinestretch}{1.5}

\renewcommand{\theenumi}{\arabic{enumi}}
\renewcommand{\labelenumi}{\arabic{enumi}.}
\renewcommand{\theenumii}{\arabic{enumi}.\arabic{enumii}}
\renewcommand{\labelenumii}{\arabic{enumi}.\arabic{enumii}.}

\newcommand{\resultfigure}[3]{%
\begin{figure}[H]
    \centering
    \includegraphics[width=0.92\textwidth]{#1}
    \caption{#2}
    \label{#3}
\end{figure}
}

\newcommand{\sourcefile}[1]{\path{#1}}

\begin{document}

\begin{titlepage}
\thispagestyle{empty}

\begin{center}
\textbf{ПРАВИТЕЛЬСТВО РОССИЙСКОЙ ФЕДЕРАЦИИ}\\
\textbf{ФГАОУ ВО НАЦИОНАЛЬНЫЙ ИССЛЕДОВАТЕЛЬСКИЙ УНИВЕРСИТЕТ}\\
\textbf{«ВЫСШАЯ ШКОЛА ЭКОНОМИКИ»}
\end{center}

\vspace{0.8cm}

\begin{center}
Факультет компьютерных наук\\
Образовательная программа «Прикладная математика и информатика»
\end{center}

\vfill

\begin{center}
\textbf{Отчет об исследовательском проекте}\\
на тему\\
\textbf{Тестирование методов оптимизации машинно-обучаемых межатомных потенциалов}
\end{center}

\vfill

\noindent
\begin{tabularx}{\textwidth}{@{}lX@{}}
\textbf{Выполнен студентом:} & \\
Группы БПМИ249 & Заренков Пётр Иванович\\
 & \textit{ФИО студента}
\end{tabularx}

\vspace{0.8cm}

\noindent
\textbf{Проверен руководителем проекта:}\\
\begin{center}
Новиков Иван Сергеевич, к. ф.-м. н.\\[-0.2cm]
\makebox[\textwidth]{\hrulefill}\\
\textit{ФИО, научная степень (если есть)}\\[0.6cm]
Доцент\\[-0.2cm]
\makebox[\textwidth]{\hrulefill}\\
\textit{Должность}\\[0.6cm]
Факультет компьютерных наук НИУ ВШЭ\\[-0.2cm]
\makebox[\textwidth]{\hrulefill}\\
\textit{Место работы (организация или департамент НИУ ВШЭ)}
\end{center}

\vfill

\begin{center}
Москва 2026
\end{center}

\end{titlepage}

\newpage
\setcounter{page}{2}

{
    \hypersetup{linkcolor=black}
    \tableofcontents
}

\newpage

\section*{Аннотация}
\addcontentsline{toc}{section}{Аннотация}
Работа посвящена сравнению методов оптимизации при обучении машинно-обучаемых межатомных потенциалов. Рассматриваются две модели: Moment Tensor Potential (MTP) и Equivariant Tensor Network (ETN), обучаемые по энергиям и силам атомных конфигураций. В качестве оптимизаторов используются native BFGS из MLIP-4, SciPy BFGS, SGD, Momentum, Adam, Adam + L-BFGS и Muon. Эксперименты проведены на системах AgPd и MoNbTaVW. По результатам расчётов квазиньютоновские методы остаются важной группой baseline-методов; для MTP на AgPd минимальную validation EPA RMSE даёт native BFGS, а SciPy BFGS остаётся близким к нему по энергетической ошибке. Для ETN native BFGS также сохраняет роль устойчивого full-batch baseline, особенно по ошибке сил и времени обучения. Среди mini-batch методов картина зависит от модели, датасета и выбранной метрики: для MTP на AgPd лучший результат даёт Adam, для ETN на AgPd Adam лидирует по energy error, а варианты Muon дают меньшую force error; на MoNbTaVW варианты Muon существенно снижают validation EPA RMSE относительно Adam, SGD и Momentum и становятся лучшими по энергетической ошибке.

\section*{Ключевые слова}
\addcontentsline{toc}{section}{Ключевые слова}
Машинно-обучаемые межатомные потенциалы; MTP; ETN; регрессия энергий и сил; BFGS; L-BFGS; SGD; Momentum; Adam; Muon; SciPy; MLIP-4; mini-batch optimization; gradient clipping; validation RMSE.

\newpage

\section{Введение}

Машинно-обучаемые межатомные потенциалы используются как быстрые приближения к расчёту потенциальной энергии атомной конфигурации и сил на атомах. Они позволяют заменить дорогостоящие прямые вычисления регрессионной моделью и тем самым расширяют возможности молекулярной динамики, оптимизации структур и скрининга материалов. Однако итоговое качество такого потенциала определяется не только архитектурой модели. Существенную роль играет и то, каким именно методом оптимизируются её параметры.

Обучение потенциала сводится к минимизации функционала, объединяющего ошибки по энергиям и силам. В этой постановке одновременно присутствуют разные масштабы компонент loss, большое число силовых наблюдений по сравнению с числом энергетических значений, плохая обусловленность поверхности ошибки и зависимость от начальной точки. Поэтому один и тот же класс потенциалов может показывать заметно разные ошибки валидации при смене оптимизатора.

В этой работе сравниваются Moment Tensor Potential (MTP) \cite{shapeev2016mtp,novikov2021mlip} и Equivariant Tensor Network (ETN) \cite{hodapp2025etn}. Для них рассматриваются native BFGS из MLIP-4, SciPy BFGS, SGD, Momentum, Adam \cite{kingma2015adam}, Adam + L-BFGS и Muon. Методы сопоставляются по validation EPA RMSE, validation forces RMSE, времени обучения, траекториям loss и разбросу по ансамблю запусков.

В основных mini-batch сериях для MTP использовалось 8000 шагов оптимизации. Для ETN на AgPd методы Adam и Muon запускались на 10000 шагов, а для ETN на MoNbTaVW на 10000 шагов запускались все mini-batch методы, поскольку при меньшем числе шагов их траектории не выходили на плато. Для BFGS-вариантов сохранялся full-batch режим с лимитом 5000 итераций.

Цель работы состоит в том, чтобы сравнить методы оптимизации при обучении MTP и ETN по качеству, времени обучения и устойчивости результатов.

Для достижения цели решаются следующие задачи:
\begin{enumerate}
    \item обучить MTP и ETN на AgPd с использованием baseline BFGS;
    \item реализовать интерфейс для обучения потенциалов через SciPy BFGS;
    \item реализовать и протестировать mini-batch оптимизаторы SGD, Momentum, Adam и Muon;
    \item провести настройку гиперпараметров для SGD и Adam;
    \item сравнить оптимизаторы на MTP и ETN;
    \item проверить перенос выводов на MoNbTaVW;
    \item проанализировать устойчивость обучения по ансамблю запусков.
\end{enumerate}

Структура работы следующая. В разделе~\ref{sec:theory} формулируется задача регрессии энергий и сил и описываются модели MTP и ETN. В разделе~\ref{sec:methods} приведены методы оптимизации. В разделе~\ref{sec:implementation} описаны вычислительный протокол и программная реализация. В разделе~\ref{sec:experiments} собраны результаты экспериментов. В разделе~\ref{sec:discussion} обсуждаются основные закономерности. Затем формулируются итоговые выводы и направления дальнейшей работы.

\newpage

\section{Теоретические основы машинно-обучаемых потенциалов}
\label{sec:theory}

\subsection{Постановка задачи регрессии энергий и сил}

Атомную конфигурацию будем записывать как набор
\[
    X = \{(\mathbf r_i, z_i)\}_{i=1}^{N},
\]
где \(\mathbf r_i \in \mathbb R^3\) -- координаты атома, \(z_i\) -- тип атома, \(N\) -- число атомов в конфигурации. Машинно-обучаемый межатомный потенциал строит приближение к потенциальной энергии \(E^{\mathrm{ref}}(X)\) и силам \(\mathbf F^{\mathrm{ref}}_{i}(X)\), взятым из эталонного расчёта или подготовленного датасета.

В локальных потенциалах полная энергия представляется как сумма атомных вкладов:
\[
    E_\theta(X) = \sum_{i=1}^{N} V_\theta(\mathcal N_i),
\]
где \(\theta\) -- параметры модели, \(\mathcal N_i\) -- локальное окружение атома \(i\) в пределах радиуса отсечения. Силы вычисляются как производные энергии по координатам:
\[
    \mathbf F_{i,\theta}(X) =
    -\frac{\partial E_\theta(X)}{\partial \mathbf r_i}.
\]

Обучение далее рассматривается как минимизация функционала ошибки на обучающей выборке \(\{X_s\}_{s=1}^{S}\):
\[
\mathcal L(\theta)
=
w_E \sum_{s=1}^{S}
\left(
E_\theta(X_s) - E_s^{\mathrm{ref}}
\right)^2
+
w_F \sum_{s=1}^{S}
\sum_{i=1}^{N_s}
\left\|
\mathbf F_{i,\theta}(X_s)
-
\mathbf F_{i,s}^{\mathrm{ref}}
\right\|^2 .
\]
В MLIP-4 значение loss и метрики ошибок вычисляются объектом \texttt{LossFunction}; в результатах дополнительно используются RMSE по полной энергии, энергии на атом (EPA RMSE) и силам.

\subsection{Moment Tensor Potential}

Moment Tensor Potential (MTP) относится к локальным межатомным потенциалам: атомная энергия в нём раскладывается по инвариантным базисным функциям локального окружения \cite{shapeev2016mtp,novikov2021mlip}. Для атома \(i\) вводится относительный радиус-вектор
\[
    \mathbf r_{ij} = \mathbf r_j - \mathbf r_i .
\]
Моментный тензор определяется формулой
\[
M_{\mu,\nu}(\mathcal N_i)
=
\sum_{j \in \mathcal N_i}
f_\mu
\left(
|\mathbf r_{ij}|, z_i, z_j
\right)
\underbrace{
\mathbf r_{ij} \otimes \cdots \otimes \mathbf r_{ij}
}_{\nu \text{ раз}},
\]
где \(\mu\) индексирует радиальную часть, \(\nu\) задаёт тензорный порядок, а \(z_i,z_j\) задают типы центрального и соседнего атомов.

Радиальная часть имеет вид
\[
f_\mu(r, z_i, z_j)
=
\sum_{k=1}^{K}
c_{\mu k z_i z_j}
\varphi_k(r),
\]
где \(\varphi_k(r)\) -- радиальные базисные функции, а \(c_{\mu k z_i z_j}\) -- обучаемые коэффициенты.

Базисные инварианты строятся как полные свёртки моментных тензоров:
\[
B_\alpha(\mathcal N_i)
=
\operatorname{contr}
\left(
M_{\mu_1,\nu_1},
M_{\mu_2,\nu_2},
\dots,
M_{\mu_m,\nu_m}
\right),
\]
где \(\operatorname{contr}\) обозначает полную свёртку тензорных индексов, дающую скалярный инвариант. Атомная энергия MTP задаётся линейной комбинацией таких инвариантов:
\[
V_\theta(\mathcal N_i)
=
\sum_{\alpha}
\xi_\alpha
B_\alpha(\mathcal N_i),
\qquad
E_\theta(X)
=
\sum_{i=1}^{N}
V_\theta(\mathcal N_i).
\]

Тензор \(M_{\mu,\nu}\) одновременно несёт радиальную и угловую информацию о локальном окружении. При \(\nu=0\) он даёт скалярный признак, при \(\nu=1\) -- векторный, при \(\nu=2\) -- тензор второго порядка. Последующие свёртки переводят эти объекты в скалярные функции, инвариантные относительно поворотов конфигурации. Параметр \texttt{level} регулирует сложность MTP: при его увеличении в модель попадают более сложные произведения и свёртки моментных тензоров.

\subsection{Equivariant Tensor Network}

Второй класс моделей в работе -- эквивариантный тензорно-сетевой потенциал ETN \cite{orus2014tnintro,stoudenmire2016tnml,hodapp2025etn}. Как и MTP, он использует локальное разложение энергии по атомным окружениям:
\begin{equation}
    E_\theta(X)
    =
    \sum_{i=1}^{N}
    V_\theta(\mathcal N_i),
\end{equation}
где \(X=\{(\mathbf r_i,z_i)\}_{i=1}^{N}\) --- атомная конфигурация, \(\mathbf r_i\) --- координаты атома, \(z_i\) --- тип атома, \(\mathcal N_i\) --- локальное окружение атома \(i\) в пределах заданного радиуса отсечения, а \(V_\theta(\mathcal N_i)\) --- атомный вклад в энергию.

Главное отличие ETN от MTP связано со способом построения признаков локального окружения. В MTP заранее задаётся набор моментных тензоров и инвариантных базисных функций, а атомная энергия записывается как их линейная комбинация. В ETN локальное окружение сначала кодируется набором эквивариантных признаков, а затем эти признаки сворачиваются в скалярную атомную энергию с помощью тензорной сети.

Для пары атомов \(i,j\) вводится относительный радиус-вектор
\begin{equation}
    \mathbf r_{ij}
    =
    \mathbf r_j - \mathbf r_i,
    \qquad
    r_{ij}
    =
    \|\mathbf r_{ij}\|,
    \qquad
    \widehat{\mathbf r}_{ij}
    =
    \frac{\mathbf r_{ij}}{r_{ij}}.
\end{equation}
Радиальная зависимость описывается набором радиальных базисных функций
\begin{equation}
    \varphi_k(r_{ij}),
    \qquad
    k=1,\ldots,K.
\end{equation}
В используемых экспериментах применялся радиальный базис \texttt{RadialBasisCinf} с параметрами \texttt{size=8}, \texttt{min\_dist=1.9}, \texttt{cutoff=5.0}. Химические типы атомов учитываются через зависимость коэффициентов модели от \(z_i\) и \(z_j\).

Эквивариантные признаки можно схематически записать в виде
\begin{equation}
    h^{(\ell)}_{i,c,m}
    =
    \sum_{j\in \mathcal N_i}
    \sum_{k}
    a^{(\ell)}_{c k}(z_i,z_j)
    \varphi_k(r_{ij})
    Y_{\ell m}(\widehat{\mathbf r}_{ij}),
\end{equation}
где \(\ell\) --- угловой порядок признака, \(m=-\ell,\ldots,\ell\) --- индекс компоненты в неприводимом представлении вращений, \(c\) --- номер канала, \(Y_{\ell m}\) --- сферические гармоники, а \(a^{(\ell)}_{c k}(z_i,z_j)\) --- обучаемые коэффициенты смешивания радиальных и химических признаков. Эта формула не фиксирует все детали реализации ETN, но отражает основной принцип: радиальная часть зависит от расстояний, а угловая часть задаёт поведение признаков при поворотах.

Ключевое свойство таких признаков состоит в эквивариантности. Если вся конфигурация поворачивается ортогональной матрицей \(Q\), то признаки порядка \(\ell\) преобразуются согласованно с представлением группы вращений:
\begin{equation}
    h^{(\ell)}_{i,c}(QX)
    =
    D^{(\ell)}(Q)
    h^{(\ell)}_{i,c}(X),
\end{equation}
где \(D^{(\ell)}(Q)\) --- матрица представления поворота \(Q\) для углового порядка \(\ell\). При \(\ell=0\) признак является скаляром и не меняется при повороте. При \(\ell=1\) признак преобразуется как вектор, а при больших \(\ell\) --- как тензор более высокого углового порядка.

Атомная энергия должна быть инвариантной величиной. Поэтому в ETN эквивариантные признаки сворачиваются так, чтобы итоговый результат имел нулевой угловой порядок:
\begin{equation}
    V_\theta(\mathcal N_i)
    =
    \operatorname{TN}_\theta
    \left(
        \{h^{(\ell)}_{i,c}\}_{\ell,c}
    \right)
    \in \mathbb R.
\end{equation}
Здесь \(\operatorname{TN}_\theta\) обозначает тензорную сеть с обучаемыми параметрами. Промежуточные признаки могут быть векторными или тензорными, но итоговая атомная энергия остаётся скаляром. Поэтому полная энергия удовлетворяет условию инвариантности:
\begin{equation}
    E_\theta(QX)
    =
    E_\theta(X).
\end{equation}

Силы вычисляются как производные энергии по координатам:
\begin{equation}
    \mathbf F_{i,\theta}(X)
    =
    -
    \frac{\partial E_\theta(X)}
    {\partial \mathbf r_i}.
\end{equation}
Из инвариантности энергии следует корректное преобразование сил при повороте конфигурации:
\begin{equation}
    \mathbf F_{i,\theta}(QX)
    =
    Q\mathbf F_{i,\theta}(X).
\end{equation}

В практической реализации ETN задаётся несколькими группами гиперпараметров. Параметр \texttt{etn\_shape} определяет ширину набора каналов для разных уровней или порядков признаков. В базовых экспериментах использовалось
\begin{equation}
    \texttt{etn\_shape}
    =
    [[1],[6],[3]].
\end{equation}
Увеличение \texttt{etn\_shape} повышает выразительность модели, но одновременно увеличивает число параметров и усложняет задачу оптимизации. Параметр \texttt{ett\_ranks} задаёт ранги тензорной сети:
\begin{equation}
    \texttt{ett\_ranks}
    =
    [[1],[1]].
\end{equation}
Малые ранги дают более компактную модель, а большие ранги увеличивают число степеней свободы. Поэтому \texttt{etn\_shape} и \texttt{ett\_ranks} влияют не только на аппроксимационную способность ETN, но и на то, насколько сложной становится задача оптимизации.

В итоге ETN можно рассматривать как модель, где физические симметрии встроены непосредственно в архитектуру. В отличие от MTP, где набор инвариантных базисных функций задан явно, ETN использует тензорно-сетевое смешивание эквивариантных признаков. Такая конструкция делает модель более гибкой, но одновременно усложняет ландшафт функции потерь. Поэтому результаты, полученные для MTP, нельзя автоматически переносить на ETN; именно это проверяется в экспериментальной части.

\section{Методы оптимизации}
\label{sec:methods}

Обучение машинно-обучаемого потенциала будем рассматривать как минимизацию функции потерь по параметрам модели. Пусть \(\theta\) обозначает вектор всех обучаемых параметров потенциала, а \(\mathcal L(\theta)\) --- функционал ошибки по энергиям и силам. Тогда задача обучения записывается как
\begin{equation}
    \theta^\ast
    =
    \arg\min_{\theta}
    \mathcal L(\theta).
\end{equation}
В работе намеренно сравниваются методы с разным объёмом информации о функции потерь. Методы первого порядка опираются только на градиент. Квазиньютоновские методы дополнительно накапливают приближение кривизны. Отдельно рассматривается Muon: в нём матричное направление обновления нормализуется с помощью ортогонализации.

\subsection{Градиентный спуск и стохастический градиентный спуск}

Самый простой метод первого порядка -- градиентный спуск. На каждой итерации параметры обновляются в направлении антиградиента:
\begin{equation}
    \theta_{t+1}
    =
    \theta_t
    -
    \eta_t
    \nabla \mathcal L(\theta_t),
\end{equation}
где \(\eta_t>0\) --- шаг обучения. Такой метод использует полный градиент по всей обучающей выборке, поэтому при фиксированных параметрах направление шага детерминировано. Недостаток состоит в стоимости одной итерации: она растёт вместе с числом конфигураций и числом атомов.

Стохастический градиентный спуск использует градиент не по всей выборке, а по мини-пакету \(\mathcal B_t\):
\begin{equation}
    g_t
    =
    \nabla
    \mathcal L_{\mathcal B_t}(\theta_t),
    \qquad
    \theta_{t+1}
    =
    \theta_t
    -
    \eta_t g_t.
\end{equation}
Такой шаг дешевле, но \(g_t\) является шумной оценкой полного градиента. Поэтому SGD вводит дополнительный компромисс: стоимость итерации уменьшается, а дисперсия направления шага растёт. Для обучения межатомных потенциалов это особенно заметно, поскольку вклад сил даёт большое число слагаемых в функции потерь, а масштабы градиентов по разным параметрам могут существенно отличаться.

На практике результат SGD сильно зависит от learning rate, размера mini-batch и ограничения нормы градиента. Слишком большой шаг может увести параметры в область больших значений loss, а слишком малый делает обучение медленным и часто приводит к затяжному плато.

\subsection{Momentum}

Momentum расширяет SGD за счёт накопления направления движения. Вместо использования только текущего градиента вводится вспомогательная переменная \(v_t\):
\begin{equation}
    v_{t+1}
    =
    \beta v_t
    +
    g_t,
\end{equation}
\begin{equation}
    \theta_{t+1}
    =
    \theta_t
    -
    \eta_t v_{t+1},
\end{equation}
где \(\beta\in[0,1)\) --- коэффициент инерции. Если последовательные градиенты направлены примерно одинаково, величина \(v_t\) накапливается, и движение в этом направлении ускоряется. Если градиенты колеблются, усреднение частично подавляет зигзагообразное движение.

Momentum полезен на плохо обусловленных поверхностях функции потерь. В таких задачах кривизна в разных направлениях может сильно отличаться: в одних направлениях функция меняется резко, в других --- медленно. Обычный SGD в этой ситуации делает много поперечных колебаний. Momentum сглаживает направление шага и может ускорять движение вдоль вытянутых долин функции потерь. При этом метод всё равно сохраняет один общий learning rate и не снимает полностью проблему разных масштабов координат.

\subsection{Adam}

Adam использует адаптивный шаг. В отличие от SGD и Momentum, он хранит две экспоненциально сглаженные статистики: первый момент градиента и второй момент покоординатного квадрата градиента. Для mini-batch градиента \(g_t\) обновления имеют вид
\begin{equation}
    m_t
    =
    \beta_1 m_{t-1}
    +
    (1-\beta_1)g_t,
\end{equation}
\begin{equation}
    v_t
    =
    \beta_2 v_{t-1}
    +
    (1-\beta_2)(g_t\odot g_t),
\end{equation}
где \(\odot\) обозначает покоординатное произведение. Поскольку в начале обучения \(m_t\) и \(v_t\) смещены к нулю, используются bias-correction оценки:
\begin{equation}
    \widehat m_t
    =
    \frac{m_t}{1-\beta_1^t},
    \qquad
    \widehat v_t
    =
    \frac{v_t}{1-\beta_2^t}.
\end{equation}
После этого параметры обновляются по правилу
\begin{equation}
    \theta_{t+1}
    =
    \theta_t
    -
    \eta_t
    \frac{\widehat m_t}
    {\sqrt{\widehat v_t}+\varepsilon}.
\end{equation}
Первый момент \(m_t\) задаёт сглаженное направление, близкое по смыслу к Momentum. Второй момент \(v_t\) нормирует шаг по координатам: если по некоторой координате градиенты систематически велики, эффективный шаг уменьшается; если градиенты малы, относительный шаг увеличивается. За счёт этого Adam частично компенсирует различие масштабов параметров и обычно устойчивее SGD при шумных mini-batch градиентах \cite{kingma2015adam}. В экспериментальной части Adam используется как основной mini-batch baseline.

\subsection{Ограничение нормы градиента}

Для mini-batch методов дополнительно применялось ограничение нормы градиента. Пусть \(C>0\) --- заданный порог clipping. Тогда вместо исходного градиента \(g_t\) используется
\begin{equation}
    g_t^{\mathrm{clip}}
    =
    g_t
    \cdot
    \min
    \left(
        1,
        \frac{C}{\|g_t\|}
    \right).
\end{equation}
Если \(\|g_t\|\leq C\), градиент не изменяется. Если \(\|g_t\|>C\), он масштабируется до нормы \(C\). Направление при этом сохраняется, а максимальный размер обновления ограничивается:
\begin{equation}
    \|\Delta \theta_t\|
    \lesssim
    \eta_t C.
\end{equation}
Для обучения потенциалов такой приём важен из-за резких изменений градиента, связанных с вкладом сил и случайным составом mini-batch. Слишком малое значение \(C\) делает обучение чрезмерно осторожным, а слишком большое почти не влияет на шаг.

\subsection{BFGS}

Методы первого порядка используют только локальный градиент. На плохо обусловленных задачах этой информации часто недостаточно: один и тот же learning rate должен одновременно подходить и для направлений с большой кривизной, и для почти плоских направлений. Квазиньютоновские методы решают эту проблему иначе: они строят приближение кривизны функции потерь и используют его при выборе направления шага.

В методе BFGS направление обновления задаётся формулой
\begin{equation}
    p_k
    =
    -
    H_k
    \nabla \mathcal L(\theta_k),
\end{equation}
где \(H_k\) --- приближение обратного гессиана. Затем выполняется шаг
\begin{equation}
    \theta_{k+1}
    =
    \theta_k
    +
    \alpha_k p_k,
\end{equation}
где \(\alpha_k\) выбирается процедурой line search или заданным правилом выбора шага.

Для обновления приближения обратного гессиана вводятся величины
\begin{equation}
    s_k
    =
    \theta_{k+1}-\theta_k,
    \qquad
    y_k
    =
    \nabla \mathcal L(\theta_{k+1})
    -
    \nabla \mathcal L(\theta_k).
\end{equation}
Обновление BFGS имеет вид
\begin{equation}
    H_{k+1}
    =
    \left(
        I
        -
        \rho_k s_k y_k^\top
    \right)
    H_k
    \left(
        I
        -
        \rho_k y_k s_k^\top
    \right)
    +
    \rho_k s_k s_k^\top,
\end{equation}
где
\begin{equation}
    \rho_k
    =
    \frac{1}{y_k^\top s_k}.
\end{equation}
BFGS не требует явного вычисления полного гессиана, но постепенно накапливает информацию о локальной кривизне через изменения градиента. Поэтому метод хорошо подходит для гладких задач с точным градиентом \cite{nocedal2006numerical}. В обучении MTP и ETN это важно: функция потерь по энергиям и силам гладко зависит от параметров потенциала, а full-batch градиент может быть вычислен средствами MLIP-4.

\subsection{L-BFGS}

Главное ограничение полного BFGS связано с хранением матрицы \(H_k\in\mathbb R^{n\times n}\), где \(n\) --- число параметров модели. Затраты памяти составляют \(O(n^2)\), что для больших моделей быстро становится существенным ограничением.

L-BFGS -- вариант BFGS с ограниченной памятью. Вместо полной матрицы \(H_k\) метод хранит только последние \(m\) пар
\begin{equation}
    (s_{k-m},y_{k-m}),\ldots,(s_{k-1},y_{k-1}).
\end{equation}
Действие приближения обратного гессиана на градиент вычисляется не через явное умножение на матрицу, а через двухпроходную рекурсию по сохранённым векторам. Поэтому память уменьшается до \(O(mn)\), где \(m\ll n\) \cite{nocedal2006numerical}.

В работе L-BFGS рассматривался в комбинированной схеме Adam + L-BFGS. Идея состоит в разделении этапов обучения: Adam сначала выводит параметры в область меньших значений функции потерь, после чего L-BFGS выполняет более точную дооптимизацию с учётом локальной кривизны. Такой подход проверялся как компромисс между устойчивостью адаптивных mini-batch методов и точностью квазиньютоновской оптимизации.

\subsection{Muon}

Muon рассматривался как дополнительный метод для параметров, которые можно представить в матричной форме. В отличие от Adam с покоординатной нормализацией шага, Muon работает с двумерными параметрами как с матрицами и нормализует само матричное направление обновления. В исходной постановке метод применяется к матричным параметрам скрытых слоёв нейронных сетей, а скалярные, векторные, входные и выходные параметры рекомендуется оптимизировать стандартным методом, например AdamW~\cite{jordan2024muon}. Связанные интерпретации матричной нормализации обновлений обсуждаются в работе~\cite{bernstein2024oldoptimizer}.

Базовый шаг Muon начинается с momentum-обновления. Пусть \(G_t\) --- градиент по матричному параметру \(W_t\in\mathbb R^{n\times m}\). Тогда сначала строится momentum-направление
\begin{equation}
    M_t
    =
    \beta M_{t-1}
    +
    G_t.
\end{equation}
Далее это матричное направление ортогонализуется:
\begin{equation}
    \widetilde M_t
    =
    \operatorname{Ortho}(M_t).
\end{equation}
После этого параметр обновляется по правилу
\begin{equation}
    W_{t+1}
    =
    W_t
    -
    \eta_t
    \widetilde M_t.
\end{equation}

Оператор \(\operatorname{Ortho}\) заменяет матрицу обновления на ближайшую полуортогональную матрицу:
\begin{equation}
    \operatorname{Ortho}(G)
    =
    \arg\min_{O}
    \left\{
        \|O-G\|_F:
        O^\top O=I
        \ \text{или}\
        OO^\top=I
    \right\}.
\end{equation}
Если сингулярное разложение матрицы \(G\) имеет вид
\begin{equation}
    G
    =
    U S V^\top,
\end{equation}
то ортогонализованное направление равно
\begin{equation}
    \operatorname{Ortho}(G)
    =
    U V^\top.
\end{equation}
Тем самым Muon сохраняет сингулярные направления обновления, но выравнивает его сингулярные значения. Это уменьшает доминирование нескольких крупных направлений и делает менее выраженные направления матричного обновления более заметными.

Прямое вычисление SVD на каждом шаге слишком дорого, поэтому в Muon ортогонализация приближённо выполняется через итерацию Newton--Schulz. Сначала матрица нормируется:
\begin{equation}
    X_0
    =
    \frac{G}{\|G\|_F+\varepsilon}.
\end{equation}
Затем выполняется несколько итераций вида
\begin{equation}
    A_t
    =
    X_t X_t^\top,
\end{equation}
\begin{equation}
    X_{t+1}
    =
    aX_t
    +
    \left(
        bA_t
        +
        cA_t^2
    \right)
    X_t,
\end{equation}
где в исходной реализации используются коэффициенты
\begin{equation}
    a=3.4445,
    \qquad
    b=-4.7750,
    \qquad
    c=2.0315.
\end{equation}
Если \(G=USV^\top\), то такая итерация действует на сингулярные значения матрицы через полином
\begin{equation}
    \varphi(x)
    =
    ax
    +
    bx^3
    +
    cx^5.
\end{equation}
После нескольких шагов сингулярные значения приближаются к единице, а матрица обновления приближается к \(UV^\top\). Поэтому Muon можно интерпретировать как momentum-метод, в котором матричное направление дополнительно приводится к полуортогональному виду.

В экспериментах Muon проверялся в двух вариантах: \texttt{Muon pad sqrt} и \texttt{Muon factorization}. Здесь важна не теоретическая демонстрация преимущества метода, а его практическое сравнение с SGD, Momentum, Adam и BFGS при обучении MTP и ETN. Результаты показывают, что качество Muon зависит от модели, датасета и метрики: на AgPd метод не превосходит Adam по validation EPA RMSE и уступает BFGS-вариантам, но для ETN даёт меньшую validation forces RMSE, чем Adam; на MoNbTaVW для ETN оба варианта Muon существенно лучше Adam, SGD и Momentum по validation EPA RMSE и дают меньшую energy error, чем BFGS-варианты, хотя по forces RMSE и времени обучения всё ещё уступают full-batch методам.

\section{Реализация вычислительного протокола}
\label{sec:implementation}

\subsection{Конфигурация экспериментов}

Расчёты выполнялись в MLIP-4. Обучающие и валидационные данные хранились в формате JSON; при запуске из них строились объекты \texttt{LossFunction}, к которым подключался текущий потенциал. Для MTP и ETN использовалась одна радиальная база:
\[
    \texttt{RadialBasisCinf(size=8, min\_dist=1.9, cutoff=5.0)}.
\]
При параллельном запуске данные и начальные параметры потенциала передавались между MPI-процессами в сериализованном виде. За счёт этого все процессы одного запуска работали с одной и той же начальной точкой и одним набором loss terms.

Вычисления запускались через Slurm на одном узле с восемью MPI-процессами. Использовались параметры одного узла \(\texttt{-N 1}\), восьми процессов \(\texttt{-n 8}\), partition \(\texttt{normal}\) и project account \(\texttt{proj\_1786}\). Отдельный запуск потенциала, соответственно, выполнялся в фиксированной восьмиядерной конфигурации.

Основные параметры моделей приведены в таблице~\ref{tab:model-config}. Для всех серий использовались пять независимых запусков с разными начальными параметрами. Номер запуска обозначался \(\texttt{pot\_num}=1,\ldots,5\), а случайная инициализация задавалась seed \(42+\texttt{pot\_num}\).

\begin{table}[H]
    \centering
    \caption{Основные параметры моделей в вычислительных экспериментах.}
    \label{tab:model-config}
    \small
    \begin{tabularx}{\textwidth}{p{0.24\textwidth}X}
        \toprule
        Серия экспериментов & Параметры \\
        \midrule
        MTP, AgPd &
        \(\texttt{species\_order}=[0,1]\);
        \(\texttt{level}=16\);
        \(\texttt{jit=True}\);
        начальная инициализация параметров из равномерного распределения на интервале \([-0.1,0.1]\). \\
        \addlinespace
        ETN, AgPd &
        \(\texttt{species\_order}=[0,1]\);
        \(\texttt{etn\_shape}=\texttt{[[1],[6],[3]]}\);
        \(\texttt{ett\_ranks}=\texttt{[[1],[1]]}\);
        \(\texttt{jit=True}\);
        начальная инициализация параметров на интервале \([-0.001,0.001]\). \\
        \addlinespace
        ETN, MoNbTaVW &
        \(\texttt{species\_order}=[0,1,2,3,4]\);
        \(\texttt{etn\_shape}=\texttt{[[1],[6],[3]]}\);
        \(\texttt{ett\_ranks}=\texttt{[[2],[2]]}\);
        \(\texttt{jit=True}\);
        начальная инициализация параметров на интервале \([-0.001,0.001]\). \\
        \bottomrule
    \end{tabularx}
\end{table}

Параметры оптимизации приведены в таблице~\ref{tab:optimization-config}. Native BFGS и SciPy BFGS использовали полный функционал ошибки и лимит \(5000\) итераций. Для SciPy BFGS дополнительно сохранялось фактическое число итераций \texttt{nit}; на ETN MoNbTaVW оно составило \(237.2 \pm 69.9\). Mini-batch методы работали с batch size 32 и постоянным learning rate. Для MTP использовалось \(8000\) шагов оптимизации. Для ETN на AgPd методы Adam и Muon запускались на \(10000\) шагов, а для ETN на MoNbTaVW на \(10000\) шагов запускались все mini-batch методы, поскольку при меньшем числе шагов их траектории ещё не выходили на плато.

\begin{table}[H]
    \centering
    \caption{Основные параметры оптимизации.}
    \label{tab:optimization-config}
    \small
    \begin{tabularx}{\textwidth}{p{0.28\textwidth}X}
        \toprule
        Метод & Параметры \\
        \midrule
        Native BFGS &
        Полный batch; максимальное число итераций \(5000\). \\
        SciPy BFGS &
        Полный batch; метод \texttt{BFGS}; \(\texttt{maxiter}=5000\); \(\texttt{gtol}=10^{-6}\); для ETN на MoNbTaVW фактически \(237.2 \pm 69.9\) итерации, \(\texttt{success}=0\). \\
        SGD &
        \(\texttt{batch\_size}=32\); \(\texttt{max\_steps}=8000\) для MTP, \(5000\) для ETN на AgPd и \(10000\) для ETN на MoNbTaVW; \(\eta=10^{-3}\); \(\texttt{ConstantLR}\); \(\texttt{clip}=1.0\); \(\texttt{full\_batch=False}\). \\
        Momentum &
        Те же общие mini-batch параметры; коэффициент инерции \(\beta=0.9\). \\
        Adam &
        Те же общие mini-batch параметры; \(\texttt{max\_steps}=10000\) для ETN на AgPd и MoNbTaVW; \(\beta_1=0.9\), \(\beta_2=0.999\), \(\varepsilon=10^{-8}\). \\
        Muon &
        Те же общие mini-batch параметры; \(\texttt{max\_steps}=10000\) для ETN на AgPd и MoNbTaVW; \(\beta=0.95\); число Newton--Schulz итераций \(5\); \(\texttt{nesterov=True}\); \(\varepsilon=10^{-7}\); варианты \(\texttt{pad\_sqrt}\) и \(\texttt{factor}\). \\
        \bottomrule
    \end{tabularx}
\end{table}

\subsection{Пользовательский цикл оптимизации}

Для mini-batch методов был реализован единый цикл оптимизации. В нём вычисление loss и градиента отделено от правила обновления параметров, поэтому один и тот же протокол можно использовать для SGD, Momentum, Adam и Muon.

\begin{itemize}
    \item \texttt{LossBatcher} собирает mini-batch функции потерь из слагаемых полного \texttt{LossFunction}. На каждой эпохе индексы перемешиваются, после чего создаётся набор batch-\texttt{LossFunction}, присоединённых к текущему потенциалу.
    \item \texttt{MlipObjective} принимает вектор параметров, записывает его в потенциал, вычисляет loss, градиент и норму градиента. Если включён clipping, градиент масштабируется до заданного порога нормы.
    \item \texttt{OptimizerState} задаёт общий интерфейс обновления параметров. Реализации \texttt{SGD}, \texttt{Momentum}, \texttt{Adam} и \texttt{Muon} получают текущие параметры, градиент и номер шага, после чего возвращают новый вектор параметров.
    \item \texttt{LearningRateSchedule} описывает изменение learning rate. В основных экспериментах применялся постоянный шаг через \texttt{ConstantLR}.
    \item \texttt{MlipTrainer} управляет циклом обучения: создаёт mini-batches, вызывает расчёт loss и градиента, применяет оптимизатор, проверяет критерий остановки по \(\texttt{gtol}\) и ограничение по числу шагов.
    \item \texttt{TrainHistory} сохраняет loss, норму градиента, learning rate, число шагов и число эпох. Эти данные затем используются при построении графиков истории обучения.
\end{itemize}

Связь компонентов mini-batch протокола показана на рисунке~\ref{fig:custom-optimization-loop}.

\begin{figure}[H]
    \centering
    \fbox{
        \begin{minipage}{0.88\textwidth}
            \centering
            \[
            \begin{array}{c}
                \text{JSON train/validation data}
                \rightarrow
                \text{\texttt{LossFunction}}
                \rightarrow
                \text{MTP/ETN potential}
                \\[0.4em]
                \downarrow
                \\[0.4em]
                \text{\texttt{LossBatcher}}
                \rightarrow
                \text{mini-batch \texttt{LossFunction}}
                \rightarrow
                \text{\texttt{MlipObjective}}
                \\[0.4em]
                \downarrow
                \\[0.4em]
                \text{\texttt{OptimizerState}}
                \rightarrow
                \text{\texttt{MlipTrainer}}
                \rightarrow
                \text{\texttt{TrainHistory}}
                \rightarrow
                \text{validation metrics}
            \end{array}
            \]
        \end{minipage}
    }
    \caption{Схема пользовательского mini-batch протокола оптимизации.}
    \label{fig:custom-optimization-loop}
\end{figure}

Исходный код экспериментов, скрипты обучения и материалы отчёта размещены в репозитории:
\url{https://github.com/PETRUESHA/Testing-optimization-methods-in-MLP}.

Для SciPy BFGS использовался отдельный адаптер \texttt{MlipWrapper}. Он приводит интерфейс MLIP-4 к формату \texttt{scipy.optimize.minimize}: метод \texttt{value(x)} возвращает значение \(\mathcal L(x)\), а метод \texttt{grad(x)} возвращает \(\nabla\mathcal L(x)\). При каждом вызове вектор \(x\) записывается в параметры потенциала; затем MLIP-4 вычисляет loss или накапливает градиент. Чтобы избежать лишних повторных вычислений, wrapper кэширует последнюю точку, значение loss и градиент.

Итоговая схема для SciPy BFGS имеет вид
\[
    \begin{array}{c}
        \text{\texttt{LossFunction} + MTP/ETN}
        \rightarrow
        \text{\texttt{MlipWrapper}}
        \\[0.3em]
        \rightarrow
        \text{\texttt{scipy.optimize.minimize}}
        \rightarrow
        \text{optimized parameters}.
    \end{array}
\]

После завершения каждого запуска вычислялись train и validation ошибки по энергиям на атом и по силам, а также время обучения. Для ансамбля из пяти запусков в итоговые таблицы заносились среднее значение и standard deviation.

\section{Эксперименты и результаты}
\label{sec:experiments}

Численные значения в этом разделе взяты из локальных CSV-таблиц с результатами обучения. Для MTP full-batch BFGS используется лимит \(5000\) итераций, а для MTP mini-batch методов --- \(8000\) шагов оптимизации. Для ETN используется конфигурация из раздела~\ref{sec:implementation}: на AgPd методы Adam и Muon обучались \(10000\) шагов, а на MoNbTaVW все mini-batch методы обучались \(10000\) шагов. Во всех таблицах EPA RMSE измеряется в eV/atom, forces RMSE -- в eV/\AA, время обучения -- в секундах. Для ансамблей указаны mean и std.

\subsection{Метрики и базовые настройки сравнения}

Основными метриками качества являются validation EPA RMSE и validation forces RMSE. Дополнительно приводятся train EPA RMSE, train forces RMSE, train time и, где это доступно, диагностические величины SciPy BFGS. Native BFGS используется как baseline, с которым далее сравниваются остальные методы.

\begin{table}[H]
\centering
\small
\caption{Базовое сравнение MTP и ETN на AgPd при native BFGS.}
\label{tab:task1-mtp-etn}
\begingroup
\setlength{\tabcolsep}{3pt}
\begin{tabular}{@{}lcccccc@{}}
\toprule
Модель & \multicolumn{2}{c}{Validation EPA RMSE, eV/atom} & \multicolumn{2}{c}{Validation forces RMSE, eV/\AA} & \multicolumn{2}{c}{Train time, s} \\
\cmidrule(lr){2-3}\cmidrule(lr){4-5}\cmidrule(lr){6-7}
 & Mean & Std & Mean & Std & Mean & Std \\
\midrule
MTP native BFGS & 0.002104 & 0.000096 & 0.047462 & 0.000910 & 2810.26 & 273.32 \\
ETN native BFGS & 0.0138745 & \(2.70491\cdot 10^{-12}\) & 0.277115 & \(9.1027\cdot 10^{-12}\) & 60.5054 & 12.6663 \\
\bottomrule
\end{tabular}
\endgroup
\end{table}

В основной конфигурации native BFGS для MTP даёт существенно меньшие ошибки, чем native BFGS для ETN, но требует заметно большего времени обучения. Это сравнение служит только общей ориентацией по двум моделям; далее оптимизаторы анализируются отдельно для MTP и ETN.

\resultfigure{fig_05_01_mtp_etn_quality_time_agpd.png}{Сравнение качества и времени обучения MTP и ETN на AgPd.}{fig:task1-quality-time}

\subsection{Настройка архитектуры ETN}

Для ETN были проверены варианты \texttt{ett\_ranks} и \texttt{etn\_shape}. Изменение rank в этом эксперименте почти не влияло на validation EPA RMSE: для \(\texttt{[[1],[1]]}\), \(\texttt{[[2],[2]]}\) и \(\texttt{[[3],[3]]}\) средняя validation EPA RMSE осталась равной 0.0138745 eV/atom. Параметр shape, напротив, оказался заметнее: вариант \(\texttt{[[1],[8],[4]]}\) дал validation EPA RMSE 0.0125834 eV/atom и validation forces RMSE 0.164119 eV/\AA, что лучше базового \(\texttt{[[1],[6],[3]]}\).

\resultfigure{fig_05_02_etn_rank_tuning_agpd.png}{Настройка \texttt{ett\_ranks} для ETN на AgPd.}{fig:task1-etn-rank}
\resultfigure{fig_05_03_etn_shape_tuning_agpd.png}{Настройка \texttt{etn\_shape} для ETN на AgPd.}{fig:task1-etn-shape}

\subsection{MTP на AgPd}
\label{subsec:mtp-results}

В этом подразделе методы оптимизации для MTP рассматриваются последовательно. BFGS-варианты работают в full-batch режиме с лимитом \(5000\) итераций. SGD, Momentum, Adam и Muon используют mini-batch режим с \(8000\) шагами оптимизации. Для каждого метода сначала приведена таблица результатов и сравнение с native BFGS, затем показан график обучения или диагностический график.

\subsubsection{Результаты для native BFGS}

\begin{table}[H]
\centering
\small
\caption{MTP на AgPd: native BFGS.}
\label{tab:mtp-native-bfgs}
\begin{tabular}{lccc}
\toprule
Метрика & Единица & Mean & Std \\
\midrule
Train EPA RMSE & eV/atom & 0.001155 & 0.000045 \\
Train forces RMSE & eV/\AA & 0.029888 & 0.000671 \\
Validation EPA RMSE & eV/atom & 0.002104 & 0.000096 \\
Validation forces RMSE & eV/\AA & 0.047462 & 0.000910 \\
Train time & s & 2810.26 & 273.32 \\
\bottomrule
\end{tabular}
\end{table}

\resultfigure{fig_05_04_mtp_bfgs_native_vs_scipy_agpd.png}{Сравнение native BFGS и SciPy BFGS для MTP на AgPd.}{fig:mtp-native-diagnostics-5000}

В MTP-серии native BFGS даёт минимальные validation EPA RMSE и validation forces RMSE. Поэтому именно этот метод используется как основной baseline для остальных MTP-оптимизаторов.

\FloatBarrier

\subsubsection{Результаты для SciPy BFGS}

\begin{table}[H]
\centering
\small
\caption{MTP на AgPd: SciPy BFGS в сравнении с native BFGS.}
\label{tab:mtp-scipy-bfgs}
\begin{tabular}{lcccc}
\toprule
\multirow{2}{*}{Метрика, единица} & \multicolumn{2}{c}{SciPy BFGS} & \multicolumn{2}{c}{Native BFGS} \\
\cmidrule(lr){2-3}\cmidrule(lr){4-5}
 & Mean & Std & Mean & Std \\
\midrule
Train EPA RMSE, eV/atom & 0.001112 & 0.000041 & 0.001155 & 0.000045 \\
Train forces RMSE, eV/\AA & 0.031007 & 0.000967 & 0.029888 & 0.000671 \\
Validation EPA RMSE, eV/atom & 0.002333 & 0.000137 & 0.002104 & 0.000096 \\
Validation forces RMSE, eV/\AA & 0.055401 & 0.000759 & 0.047462 & 0.000910 \\
Train time, s & 2894.70 & 69.03 & 2810.26 & 273.32 \\
\bottomrule
\end{tabular}
\end{table}

\resultfigure{fig_05_05_mtp_bfgs_diagnostics_agpd.png}{Диагностика BFGS-вариантов для MTP на AgPd.}{fig:mtp-scipy-bfgs-diagnostics-5000}

SciPy BFGS близок к native BFGS по энергетической ошибке, но уступает ему по validation EPA RMSE и validation forces RMSE. Среднее число итераций SciPy BFGS составило 4873.6; success rate равен 0, что в данной серии соответствует остановке по лимиту итераций во всех запусках.

\FloatBarrier

\subsubsection{Результаты для SGD}

\begin{table}[H]
\centering
\small
\caption{MTP на AgPd: SGD в сравнении с native BFGS.}
\label{tab:mtp-sgd}
\begin{tabular}{lcccc}
\toprule
\multirow{2}{*}{Метрика, единица} & \multicolumn{2}{c}{SGD} & \multicolumn{2}{c}{Native BFGS} \\
\cmidrule(lr){2-3}\cmidrule(lr){4-5}
 & Mean & Std & Mean & Std \\
\midrule
Train EPA RMSE, eV/atom & 0.080615 & 0.022432 & 0.001155 & 0.000045 \\
Train forces RMSE, eV/\AA & 0.984533 & 0.116070 & 0.029888 & 0.000671 \\
Validation EPA RMSE, eV/atom & 0.093689 & 0.026609 & 0.002104 & 0.000096 \\
Validation forces RMSE, eV/\AA & 1.424873 & 0.146451 & 0.047462 & 0.000910 \\
Train time, s & 1958.05 & 69.93 & 2810.26 & 273.32 \\
\bottomrule
\end{tabular}
\end{table}

\resultfigure{fig_05_08_mtp_sgd_loss_agpd.png}{Сглаженная история loss для MTP SGD: mean \(\pm\) std по ансамблю.}{fig:mtp-sgd-band-8000}

SGD уступает Adam, Momentum и Muon pad sqrt по validation EPA RMSE и имеет наибольшую force error среди mini-batch методов. В выбранной конфигурации это показывает чувствительность SGD к масштабу шага и шуму mini-batch градиента.

\FloatBarrier

\subsubsection{Результаты для Momentum}

\begin{table}[H]
\centering
\small
\caption{MTP на AgPd: Momentum в сравнении с native BFGS.}
\label{tab:mtp-momentum}
\begin{tabular}{lcccc}
\toprule
\multirow{2}{*}{Метрика, единица} & \multicolumn{2}{c}{Momentum} & \multicolumn{2}{c}{Native BFGS} \\
\cmidrule(lr){2-3}\cmidrule(lr){4-5}
 & Mean & Std & Mean & Std \\
\midrule
Train EPA RMSE, eV/atom & 0.042911 & 0.008859 & 0.001155 & 0.000045 \\
Train forces RMSE, eV/\AA & 0.774203 & 0.083373 & 0.029888 & 0.000671 \\
Validation EPA RMSE, eV/atom & 0.050901 & 0.015717 & 0.002104 & 0.000096 \\
Validation forces RMSE, eV/\AA & 1.148895 & 0.096406 & 0.047462 & 0.000910 \\
Train time, s & 1924.49 & 22.12 & 2810.26 & 273.32 \\
\bottomrule
\end{tabular}
\end{table}

\resultfigure{fig_05_07_mtp_momentum_loss_agpd.png}{Сглаженная история loss для MTP Momentum: mean \(\pm\) std по ансамблю.}{fig:mtp-momentum-band-8000}

Momentum занимает второе место среди pure mini-batch методов MTP по validation EPA RMSE, но ошибка по силам остаётся большой. Поэтому улучшение energy error здесь не означает сопоставимого улучшения forces RMSE.

\FloatBarrier

\subsubsection{Результаты для Adam}

\begin{table}[H]
\centering
\small
\caption{MTP на AgPd: Adam в сравнении с native BFGS.}
\label{tab:mtp-adam}
\begin{tabular}{lcccc}
\toprule
\multirow{2}{*}{Метрика, единица} & \multicolumn{2}{c}{Adam} & \multicolumn{2}{c}{Native BFGS} \\
\cmidrule(lr){2-3}\cmidrule(lr){4-5}
 & Mean & Std & Mean & Std \\
\midrule
Train EPA RMSE, eV/atom & 0.027625 & 0.007792 & 0.001155 & 0.000045 \\
Train forces RMSE, eV/\AA & 0.201813 & 0.047051 & 0.029888 & 0.000671 \\
Validation EPA RMSE, eV/atom & 0.029321 & 0.007448 & 0.002104 & 0.000096 \\
Validation forces RMSE, eV/\AA & 0.338570 & 0.074599 & 0.047462 & 0.000910 \\
Train time, s & 1913.95 & 15.62 & 2810.26 & 273.32 \\
\bottomrule
\end{tabular}
\end{table}

\resultfigure{fig_05_06_mtp_adam_loss_agpd.png}{Сглаженная история loss для MTP Adam: mean \(\pm\) std по ансамблю.}{fig:mtp-adam-band-8000}

Среди mini-batch методов для MTP лучший validation EPA RMSE даёт Adam. При этом до BFGS-вариантов по energy error и force error остаётся большой разрыв, хотя время обучения у Adam меньше.

\FloatBarrier

\subsubsection{Результаты для Muon pad sqrt}

\begin{table}[H]
\centering
\small
\caption{MTP на AgPd: Muon pad sqrt в сравнении с native BFGS.}
\label{tab:mtp-muon-pad}
\begin{tabular}{lcccc}
\toprule
\multirow{2}{*}{Метрика, единица} & \multicolumn{2}{c}{Muon pad sqrt} & \multicolumn{2}{c}{Native BFGS} \\
\cmidrule(lr){2-3}\cmidrule(lr){4-5}
 & Mean & Std & Mean & Std \\
\midrule
Train EPA RMSE, eV/atom & 0.080235 & 0.032057 & 0.001155 & 0.000045 \\
Train forces RMSE, eV/\AA & 0.313185 & 0.052386 & 0.029888 & 0.000671 \\
Validation EPA RMSE, eV/atom & 0.080185 & 0.025868 & 0.002104 & 0.000096 \\
Validation forces RMSE, eV/\AA & 0.503803 & 0.078258 & 0.047462 & 0.000910 \\
Train time, s & 2052.81 & 237.96 & 2810.26 & 273.32 \\
\bottomrule
\end{tabular}
\end{table}

\resultfigure{fig_05_09_mtp_muon_pad_sqrt_loss_agpd.png}{Сглаженная история loss для MTP Muon pad sqrt: mean \(\pm\) std по ансамблю.}{fig:mtp-muon-pad-band-8000}

Muon pad sqrt лучше Muon factorization по validation EPA RMSE и validation forces RMSE, но Adam не превосходит. По force error он лучше Momentum и SGD, однако по energy error уступает Adam и Momentum.

\FloatBarrier

\subsubsection{Результаты для Muon factorization}

\begin{table}[H]
\centering
\small
\caption{MTP на AgPd: Muon factorization в сравнении с native BFGS.}
\label{tab:mtp-muon-factor}
\begin{tabular}{lcccc}
\toprule
\multirow{2}{*}{Метрика, единица} & \multicolumn{2}{c}{Muon factorization} & \multicolumn{2}{c}{Native BFGS} \\
\cmidrule(lr){2-3}\cmidrule(lr){4-5}
 & Mean & Std & Mean & Std \\
\midrule
Train EPA RMSE, eV/atom & 0.117925 & 0.039347 & 0.001155 & 0.000045 \\
Train forces RMSE, eV/\AA & 0.385825 & 0.037469 & 0.029888 & 0.000671 \\
Validation EPA RMSE, eV/atom & 0.113416 & 0.045063 & 0.002104 & 0.000096 \\
Validation forces RMSE, eV/\AA & 0.598674 & 0.054244 & 0.047462 & 0.000910 \\
Train time, s & 1913.01 & 12.00 & 2810.26 & 273.32 \\
\bottomrule
\end{tabular}
\end{table}

\resultfigure{fig_05_10_mtp_muon_factorization_loss_agpd.png}{Сглаженная история loss для MTP Muon factorization: mean \(\pm\) std по ансамблю.}{fig:mtp-muon-factor-band-8000}

Muon factorization уступает Muon pad sqrt по validation EPA RMSE и validation forces RMSE. Для MTP на AgPd факторизованный вариант, таким образом, не даёт преимущества.

\FloatBarrier

\subsubsection{Сводное сравнение всех методов для MTP}

\begin{table}[H]
\centering
\small
\caption{MTP на AgPd: сводное сравнение оптимизаторов. Единицы указаны в заголовках столбцов.}
\label{tab:mtp-summary}
\begingroup
\setlength{\tabcolsep}{3pt}
\begin{tabular}{@{}lccc@{}}
\toprule
Метод & Validation EPA RMSE, eV/atom & Validation forces RMSE, eV/\AA & Train time, s \\
\midrule
Native BFGS & \(0.002104 \pm 0.000096\) & \(0.047462 \pm 0.000910\) & \(2810.26 \pm 273.32\) \\
SciPy BFGS & \(0.002333 \pm 0.000137\) & \(0.055401 \pm 0.000759\) & \(2894.70 \pm 69.03\) \\
SGD & \(0.093689 \pm 0.026609\) & \(1.424873 \pm 0.146451\) & \(1958.05 \pm 69.93\) \\
Momentum & \(0.050901 \pm 0.015717\) & \(1.148895 \pm 0.096406\) & \(1924.49 \pm 22.12\) \\
Adam & \(0.029321 \pm 0.007448\) & \(0.338570 \pm 0.074599\) & \(1913.95 \pm 15.62\) \\
Muon pad sqrt & \(0.080185 \pm 0.025868\) & \(0.503803 \pm 0.078258\) & \(2052.81 \pm 237.96\) \\
Muon factorization & \(0.113416 \pm 0.045063\) & \(0.598674 \pm 0.054244\) & \(1913.01 \pm 12.00\) \\
\bottomrule
\end{tabular}
\endgroup
\end{table}

\resultfigure{fig_05_11_mtp_validation_errors_agpd.png}{Validation errors для MTP на AgPd.}{fig:mtp-validation-8000}
\resultfigure{fig_05_12_mtp_time_quality_agpd.png}{Качество и время обучения MTP на AgPd.}{fig:mtp-time-quality-8000}
\resultfigure{fig_05_13_mtp_ensemble_spread_agpd.png}{Разброс validation RMSE по ансамблю для MTP на AgPd.}{fig:mtp-ensemble-8000}
\resultfigure{fig_05_14_mtp_loss_overview_agpd.png}{Обзор сглаженных loss histories для MTP mini-batch оптимизаторов.}{fig:mtp-losses-overview-8000}

Сводное сравнение MTP показывает, что BFGS-варианты остаются сильнейшей группой методов. Native BFGS немного лучше SciPy BFGS по validation EPA RMSE и заметнее лучше по forces RMSE. Среди mini-batch методов минимальную validation EPA RMSE даёт Adam. Далее по energy error идёт Momentum; Muon pad sqrt занимает промежуточное положение, а SGD и Muon factorization уступают по одной или нескольким основным метрикам.

\FloatBarrier

\subsection{ETN на AgPd}
\label{subsec:etn-agpd-results}

Для ETN на AgPd native BFGS даёт лучшие validation EPA RMSE и validation forces RMSE. В отличие от MTP, SciPy BFGS не даёт устойчивого улучшения относительно native BFGS. Для Adam и Muon использовалось \(10000\) шагов оптимизации, поскольку при меньшем числе шагов их loss histories ещё не выходили на плато.

\subsubsection{Результаты для native BFGS}

\begin{table}[H]
\centering
\small
\caption{ETN на AgPd: native BFGS.}
\label{tab:etn-agpd-native-bfgs}
\begin{tabular}{lccc}
\toprule
Метрика & Единица & Mean & Std \\
\midrule
Validation EPA RMSE & eV/atom & 0.0138745 & \(2.70491\cdot 10^{-12}\) \\
Validation forces RMSE & eV/\AA & 0.277115 & \(9.1027\cdot 10^{-12}\) \\
Train time & s & 60.5054 & 12.6663 \\
\bottomrule
\end{tabular}
\end{table}

Для full-batch native BFGS отдельная mini-batch история loss не строилась. Этот метод используется как основной baseline для ETN на AgPd, поскольку даёт минимальные validation EPA RMSE и validation forces RMSE среди рассмотренных оптимизаторов.

\FloatBarrier

\subsubsection{Результаты для SciPy BFGS}

\begin{table}[H]
\centering
\small
\caption{ETN на AgPd: SciPy BFGS в сравнении с native BFGS.}
\label{tab:etn-agpd-scipy-bfgs}
\begin{tabular}{lcccc}
\toprule
\multirow{2}{*}{Метрика, единица} & \multicolumn{2}{c}{SciPy BFGS} & \multicolumn{2}{c}{Native BFGS} \\
\cmidrule(lr){2-3}\cmidrule(lr){4-5}
 & Mean & Std & Mean & Std \\
\midrule
Validation EPA RMSE, eV/atom & 0.0332907 & 0.0432011 & 0.0138745 & \(2.70491\cdot 10^{-12}\) \\
Validation forces RMSE, eV/\AA & 0.401276 & 0.277299 & 0.277115 & \(9.1027\cdot 10^{-12}\) \\
Train time, s & 81.4725 & 23.7965 & 60.5054 & 12.6663 \\
\bottomrule
\end{tabular}
\end{table}

\resultfigure{fig_05_15_etn_bfgs_native_vs_scipy_agpd.png}{Диагностическое сравнение native BFGS и SciPy BFGS для ETN на AgPd.}{fig:etn-agpd-bfgs-old}

SciPy BFGS уступает native BFGS по средним validation-метрикам и имеет заметный разброс по ансамблю. В отличие от MTP, внешний BFGS-интерфейс для ETN на AgPd не улучшает качество.

\FloatBarrier

\subsubsection{Результаты для SGD}

\begin{table}[H]
\centering
\small
\caption{ETN на AgPd: SGD в сравнении с native BFGS.}
\label{tab:etn-agpd-sgd}
\begin{tabular}{lcccc}
\toprule
\multirow{2}{*}{Метрика, единица} & \multicolumn{2}{c}{SGD} & \multicolumn{2}{c}{Native BFGS} \\
\cmidrule(lr){2-3}\cmidrule(lr){4-5}
 & Mean & Std & Mean & Std \\
\midrule
Validation EPA RMSE, eV/atom & 0.172721 & 0.001135 & 0.0138745 & \(2.70491\cdot 10^{-12}\) \\
Validation forces RMSE, eV/\AA & 1.73016 & \(6.61807\cdot 10^{-10}\) & 0.277115 & \(9.1027\cdot 10^{-12}\) \\
Train time, s & 1061.03 & 1.36 & 60.5054 & 12.6663 \\
\bottomrule
\end{tabular}
\end{table}

\resultfigure{fig_05_16_etn_sgd_loss_agpd.png}{Сглаженная история loss для ETN SGD на AgPd.}{fig:agpd-sgd-band-5000}

SGD на ETN даёт высокие validation errors; в этой постановке он заметно уступает остальным baseline-методам.

\FloatBarrier

\subsubsection{Результаты для Momentum}

\begin{table}[H]
\centering
\small
\caption{ETN на AgPd: Momentum в сравнении с native BFGS.}
\label{tab:etn-agpd-momentum}
\begin{tabular}{lcccc}
\toprule
\multirow{2}{*}{Метрика, единица} & \multicolumn{2}{c}{Momentum} & \multicolumn{2}{c}{Native BFGS} \\
\cmidrule(lr){2-3}\cmidrule(lr){4-5}
 & Mean & Std & Mean & Std \\
\midrule
Validation EPA RMSE, eV/atom & 0.172698 & 0.001643 & 0.0138745 & \(2.70491\cdot 10^{-12}\) \\
Validation forces RMSE, eV/\AA & 1.73016 & \(6.63209\cdot 10^{-10}\) & 0.277115 & \(9.1027\cdot 10^{-12}\) \\
Train time, s & 1063.97 & 4.44 & 60.5054 & 12.6663 \\
\bottomrule
\end{tabular}
\end{table}

\resultfigure{fig_05_17_etn_momentum_loss_agpd.png}{Сглаженная история loss для ETN Momentum на AgPd.}{fig:agpd-momentum-band-5000}

Momentum даёт близкое к SGD значение validation EPA RMSE и не улучшает force error.

\FloatBarrier

\subsubsection{Результаты для Adam}

\begin{table}[H]
\centering
\small
\caption{ETN на AgPd: Adam в сравнении с native BFGS.}
\label{tab:etn-agpd-adam}
\begin{tabular}{lcccc}
\toprule
\multirow{2}{*}{Метрика, единица} & \multicolumn{2}{c}{Adam} & \multicolumn{2}{c}{Native BFGS} \\
\cmidrule(lr){2-3}\cmidrule(lr){4-5}
 & Mean & Std & Mean & Std \\
\midrule
Validation EPA RMSE, eV/atom & 0.022675 & 0.000710 & 0.0138745 & \(2.70491\cdot 10^{-12}\) \\
Validation forces RMSE, eV/\AA & 0.870676 & 0.000541 & 0.277115 & \(9.1027\cdot 10^{-12}\) \\
Train time, s & 2034.93 & 9.39 & 60.5054 & 12.6663 \\
\bottomrule
\end{tabular}
\end{table}

\resultfigure{fig_05_18_etn_adam_loss_agpd.png}{Сглаженная история loss для ETN Adam на AgPd.}{fig:agpd-adam-band-10000}

Для ETN на AgPd Adam даёт лучшую validation EPA RMSE среди mini-batch методов. По force error он уступает вариантам Muon, но остаётся существенно лучше SGD и Momentum.

\FloatBarrier

\subsubsection{Результаты для Muon pad sqrt}

\begin{table}[H]
\centering
\small
\caption{ETN на AgPd: Muon pad sqrt в сравнении с native BFGS.}
\label{tab:etn-agpd-muon-pad}
\begin{tabular}{lcccc}
\toprule
\multirow{2}{*}{Метрика, единица} & \multicolumn{2}{c}{Muon pad sqrt} & \multicolumn{2}{c}{Native BFGS} \\
\cmidrule(lr){2-3}\cmidrule(lr){4-5}
 & Mean & Std & Mean & Std \\
\midrule
Validation EPA RMSE, eV/atom & 0.023418 & 0.000600 & 0.0138745 & \(2.70491\cdot 10^{-12}\) \\
Validation forces RMSE, eV/\AA & 0.856381 & 0.005138 & 0.277115 & \(9.1027\cdot 10^{-12}\) \\
Train time, s & 2168.03 & 265.08 & 60.5054 & 12.6663 \\
\bottomrule
\end{tabular}
\end{table}

\resultfigure{fig_05_19_etn_muon_pad_sqrt_loss_agpd.png}{Сглаженная история loss для ETN Muon pad sqrt на AgPd.}{fig:agpd-muon-pad-band-10000}

Muon pad sqrt близок к Adam по validation EPA RMSE и лучше Adam по validation forces RMSE. Относительно native BFGS он всё равно даёт более высокую ошибку и требует существенно большего времени.

\FloatBarrier

\subsubsection{Результаты для Muon factorization}

\begin{table}[H]
\centering
\small
\caption{ETN на AgPd: Muon factorization в сравнении с native BFGS.}
\label{tab:etn-agpd-muon-factor}
\begin{tabular}{lcccc}
\toprule
\multirow{2}{*}{Метрика, единица} & \multicolumn{2}{c}{Muon factorization} & \multicolumn{2}{c}{Native BFGS} \\
\cmidrule(lr){2-3}\cmidrule(lr){4-5}
 & Mean & Std & Mean & Std \\
\midrule
Validation EPA RMSE, eV/atom & 0.023494 & 0.001181 & 0.0138745 & \(2.70491\cdot 10^{-12}\) \\
Validation forces RMSE, eV/\AA & 0.850381 & 0.005353 & 0.277115 & \(9.1027\cdot 10^{-12}\) \\
Train time, s & 2037.94 & 7.76 & 60.5054 & 12.6663 \\
\bottomrule
\end{tabular}
\end{table}

\resultfigure{fig_05_20_etn_muon_factorization_loss_agpd.png}{Сглаженная история loss для ETN Muon factorization на AgPd.}{fig:agpd-muon-factor-band-10000}

Muon factorization близок к Muon pad sqrt по validation EPA RMSE и немного лучше по validation forces RMSE. Среди mini-batch методов он даёт минимальную force error, но по energy error уступает Adam.

\FloatBarrier

\subsubsection{Сводное сравнение всех методов для ETN на AgPd}

\begin{table}[H]
\centering
\small
\caption{ETN на AgPd: сводное сравнение оптимизаторов.}
\label{tab:etn-agpd-summary}
\begingroup
\setlength{\tabcolsep}{2.5pt}
\begin{tabular}{@{}lccc@{}}
\toprule
Метод & Validation EPA RMSE, eV/atom & Validation forces RMSE, eV/\AA & Train time, s \\
\midrule
Native BFGS & \(0.0138745 \pm 2.70491\cdot 10^{-12}\) & \(0.277115 \pm 9.1027\cdot 10^{-12}\) & \(60.5054 \pm 12.6663\) \\
SciPy BFGS & \(0.033291 \pm 0.043201\) & \(0.401276 \pm 0.277299\) & \(81.47 \pm 23.80\) \\
SGD & \(0.172721 \pm 0.001135\) & \(1.73016 \pm 6.61807\cdot 10^{-10}\) & \(1061.03 \pm 1.36\) \\
Momentum & \(0.172698 \pm 0.001643\) & \(1.73016 \pm 6.63209\cdot 10^{-10}\) & \(1063.97 \pm 4.44\) \\
Adam & \(0.022675 \pm 0.000710\) & \(0.870676 \pm 0.000541\) & \(2034.93 \pm 9.39\) \\
Muon pad sqrt & \(0.023418 \pm 0.000600\) & \(0.856381 \pm 0.005138\) & \(2168.03 \pm 265.08\) \\
Muon factorization & \(0.023494 \pm 0.001181\) & \(0.850381 \pm 0.005353\) & \(2037.94 \pm 7.76\) \\
\bottomrule
\end{tabular}
\endgroup
\end{table}

\resultfigure{fig_05_21_etn_validation_errors_agpd.png}{Validation errors для ETN на AgPd.}{fig:agpd-validation-10000}
\resultfigure{fig_05_22_etn_loss_overview_agpd.png}{Обзор сглаженных loss histories для ETN на AgPd.}{fig:agpd-losses-overview-10000}
\resultfigure{fig_05_23_etn_time_quality_agpd.png}{Качество и время обучения ETN на AgPd.}{fig:agpd-time-quality-10000}
\resultfigure{fig_05_24_etn_ensemble_spread_agpd.png}{Разброс ETN-результатов по ансамблю на AgPd.}{fig:agpd-ensemble-10000}

Сводное сравнение показывает, что native BFGS остаётся лучшим методом по обеим validation-метрикам. Среди mini-batch методов Adam даёт минимальную validation EPA RMSE, а варианты Muon близки к нему по energy error и немного лучше по force error. SGD и Momentum заметно уступают остальным методам.

\FloatBarrier

\subsection{ETN на MoNbTaVW}
\label{subsec:etn-monbtavw-results}

MoNbTaVW используется как проверка переносимости выводов на другой многокомпонентный датасет. В этой серии все mini-batch методы обучались \(10000\) шагов. Native BFGS и SciPy BFGS остаются сильными full-batch baseline, но по validation EPA RMSE после увеличения числа шагов лучшими оказываются варианты Muon.

\subsubsection{Результаты для native BFGS}

\begin{table}[H]
\centering
\small
\caption{ETN на MoNbTaVW: native BFGS.}
\label{tab:etn-monbtavw-native-bfgs}
\begin{tabular}{lccc}
\toprule
Метрика & Единица & Mean & Std \\
\midrule
Validation EPA RMSE & eV/atom & 0.476876 & \(2.33759\cdot 10^{-10}\) \\
Validation forces RMSE & eV/\AA & 0.523471 & \(8.20466\cdot 10^{-11}\) \\
Train time & s & 145.110 & 41.9128 \\
\bottomrule
\end{tabular}
\end{table}

Для full-batch native BFGS отдельная mini-batch история loss не строилась. На MoNbTaVW этот метод служит baseline для остальных оптимизаторов и даёт минимальную validation forces RMSE, но по validation EPA RMSE уступает вариантам Muon после \(10000\) mini-batch шагов.

\FloatBarrier

\subsubsection{Результаты для SciPy BFGS}

\begin{table}[H]
\centering
\small
\caption{ETN на MoNbTaVW: SciPy BFGS в сравнении с native BFGS.}
\label{tab:etn-monbtavw-scipy-bfgs}
\begin{tabular}{lcccc}
\toprule
\multirow{2}{*}{Метрика, единица} & \multicolumn{2}{c}{SciPy BFGS} & \multicolumn{2}{c}{Native BFGS} \\
\cmidrule(lr){2-3}\cmidrule(lr){4-5}
 & Mean & Std & Mean & Std \\
\midrule
Validation EPA RMSE, eV/atom & 0.499952 & \(4.27592\cdot 10^{-10}\) & 0.476876 & \(2.33759\cdot 10^{-10}\) \\
Validation forces RMSE, eV/\AA & 0.529053 & \(1.30138\cdot 10^{-9}\) & 0.523471 & \(8.20466\cdot 10^{-11}\) \\
Train time, s & 240.921 & 63.9803 & 145.110 & 41.9128 \\
SciPy iterations, count & 237.2 & 69.9085 & -- & -- \\
SciPy success, flag & 0 & 0 & -- & -- \\
\bottomrule
\end{tabular}
\end{table}

\resultfigure{fig_05_25_etn_bfgs_diagnostics_monbtavw.png}{Диагностическое сравнение native BFGS и SciPy BFGS для ETN на MoNbTaVW.}{fig:task6-bfgs-diagnostics}

SciPy BFGS близок к native BFGS по energy error и force error, но по обеим validation-метрикам немного уступает native BFGS. Среднее время SciPy BFGS выше: \(240.921\) s против \(145.110\) s у native BFGS. При этом среднее число итераций SciPy BFGS равно \(237.2\), а флаг \texttt{success} равен нулю во всех запусках; в данной серии это не связано с плохим качеством решения, поскольку итоговые ошибки близки к native BFGS. По validation EPA RMSE оба BFGS-варианта уступают Muon.

\FloatBarrier

\subsubsection{Результаты для SGD}

\begin{table}[H]
\centering
\small
\caption{ETN на MoNbTaVW: SGD в сравнении с native BFGS.}
\label{tab:etn-monbtavw-sgd}
\begin{tabular}{lcccc}
\toprule
\multirow{2}{*}{Метрика, единица} & \multicolumn{2}{c}{SGD} & \multicolumn{2}{c}{Native BFGS} \\
\cmidrule(lr){2-3}\cmidrule(lr){4-5}
 & Mean & Std & Mean & Std \\
\midrule
Validation EPA RMSE, eV/atom & 1.72847 & 0.000501 & 0.476876 & \(2.33759\cdot 10^{-10}\) \\
Validation forces RMSE, eV/\AA & 0.887349 & \(2.30430\cdot 10^{-10}\) & 0.523471 & \(8.20466\cdot 10^{-11}\) \\
Train time, s & 5446.38 & 108.408 & 145.110 & 41.9128 \\
\bottomrule
\end{tabular}
\end{table}

\resultfigure{fig_05_26_etn_sgd_loss_monbtavw.png}{Сглаженная история loss для ETN SGD на MoNbTaVW.}{fig:monbtavw-sgd-band-10000}

SGD остаётся слабым mini-batch методом для ETN на MoNbTaVW. Увеличение числа шагов уменьшило validation EPA RMSE по сравнению с более короткими запусками, но метод всё ещё заметно уступает Adam, Muon и BFGS-вариантам по energy error, а также проигрывает BFGS по force error и времени.

\FloatBarrier

\subsubsection{Результаты для Momentum}

\begin{table}[H]
\centering
\small
\caption{ETN на MoNbTaVW: Momentum в сравнении с native BFGS.}
\label{tab:etn-monbtavw-momentum}
\begin{tabular}{lcccc}
\toprule
\multirow{2}{*}{Метрика, единица} & \multicolumn{2}{c}{Momentum} & \multicolumn{2}{c}{Native BFGS} \\
\cmidrule(lr){2-3}\cmidrule(lr){4-5}
 & Mean & Std & Mean & Std \\
\midrule
Validation EPA RMSE, eV/atom & 1.73013 & 0.000473 & 0.476876 & \(2.33759\cdot 10^{-10}\) \\
Validation forces RMSE, eV/\AA & 0.887349 & \(2.30449\cdot 10^{-10}\) & 0.523471 & \(8.20466\cdot 10^{-11}\) \\
Train time, s & 5425.00 & 50.5405 & 145.110 & 41.9128 \\
\bottomrule
\end{tabular}
\end{table}

\resultfigure{fig_05_27_etn_momentum_loss_monbtavw.png}{Сглаженная история loss для ETN Momentum на MoNbTaVW.}{fig:monbtavw-momentum-band-10000}

Momentum ведёт себя близко к SGD: energy error остаётся существенно выше, чем у BFGS, Adam и Muon. В данной конфигурации добавление momentum не устраняет слабость метода по validation EPA RMSE.

\FloatBarrier

\subsubsection{Результаты для Adam}

\begin{table}[H]
\centering
\small
\caption{ETN на MoNbTaVW: Adam в сравнении с native BFGS.}
\label{tab:etn-monbtavw-adam}
\begin{tabular}{lcccc}
\toprule
\multirow{2}{*}{Метрика, единица} & \multicolumn{2}{c}{Adam} & \multicolumn{2}{c}{Native BFGS} \\
\cmidrule(lr){2-3}\cmidrule(lr){4-5}
 & Mean & Std & Mean & Std \\
\midrule
Validation EPA RMSE, eV/atom & 0.441821 & 0.001756 & 0.476876 & \(2.33759\cdot 10^{-10}\) \\
Validation forces RMSE, eV/\AA & 0.678150 & 0.000513 & 0.523471 & \(8.20466\cdot 10^{-11}\) \\
Train time, s & 5344.02 & 28.9838 & 145.110 & 41.9128 \\
\bottomrule
\end{tabular}
\end{table}

\resultfigure{fig_05_28_etn_adam_loss_monbtavw.png}{Сглаженная история loss для ETN Adam на MoNbTaVW.}{fig:monbtavw-adam-band-10000}

Adam существенно лучше SGD и Momentum и после \(10000\) шагов даёт меньшую validation EPA RMSE, чем native BFGS и SciPy BFGS. При этом он уступает вариантам Muon по energy error и force error, а по validation forces RMSE остаётся хуже BFGS-вариантов.

\FloatBarrier

\subsubsection{Результаты для Muon pad sqrt}

\begin{table}[H]
\centering
\small
\caption{ETN на MoNbTaVW: Muon pad sqrt в сравнении с native BFGS.}
\label{tab:etn-monbtavw-muon-pad}
\begin{tabular}{lcccc}
\toprule
\multirow{2}{*}{Метрика, единица} & \multicolumn{2}{c}{Muon pad sqrt} & \multicolumn{2}{c}{Native BFGS} \\
\cmidrule(lr){2-3}\cmidrule(lr){4-5}
 & Mean & Std & Mean & Std \\
\midrule
Validation EPA RMSE, eV/atom & 0.428698 & 0.004227 & 0.476876 & \(2.33759\cdot 10^{-10}\) \\
Validation forces RMSE, eV/\AA & 0.600915 & 0.003125 & 0.523471 & \(8.20466\cdot 10^{-11}\) \\
Train time, s & 5392.03 & 106.256 & 145.110 & 41.9128 \\
\bottomrule
\end{tabular}
\end{table}

\resultfigure{fig_05_29_etn_muon_pad_sqrt_loss_monbtavw.png}{Сглаженная история loss для ETN Muon pad sqrt на MoNbTaVW.}{fig:monbtavw-muon-pad-band-10000}

Muon pad sqrt существенно превосходит Adam, SGD и Momentum по validation EPA RMSE и становится лучше BFGS-вариантов по energy error. При этом force error остаётся хуже, чем у BFGS-вариантов, а время обучения значительно выше.

\FloatBarrier

\subsubsection{Результаты для Muon factorization}

\begin{table}[H]
\centering
\small
\caption{ETN на MoNbTaVW: Muon factorization в сравнении с native BFGS.}
\label{tab:etn-monbtavw-muon-factor}
\begin{tabular}{lcccc}
\toprule
\multirow{2}{*}{Метрика, единица} & \multicolumn{2}{c}{Muon factorization} & \multicolumn{2}{c}{Native BFGS} \\
\cmidrule(lr){2-3}\cmidrule(lr){4-5}
 & Mean & Std & Mean & Std \\
\midrule
Validation EPA RMSE, eV/atom & 0.428698 & 0.004227 & 0.476876 & \(2.33759\cdot 10^{-10}\) \\
Validation forces RMSE, eV/\AA & 0.600915 & 0.003125 & 0.523471 & \(8.20466\cdot 10^{-11}\) \\
Train time, s & 5423.30 & 63.3786 & 145.110 & 41.9128 \\
\bottomrule
\end{tabular}
\end{table}

\resultfigure{fig_05_30_etn_muon_factorization_loss_monbtavw.png}{Сглаженная история loss для ETN Muon factorization на MoNbTaVW.}{fig:monbtavw-muon-factor-band-10000}

Muon factorization даёт те же средние validation EPA RMSE и validation forces RMSE, что и Muon pad sqrt, но немного отличается по времени обучения. В этой серии оба варианта Muon приводят к одному качественному выводу: матричная нормализация обновлений особенно полезна для ETN на MoNbTaVW по энергетической ошибке.

\FloatBarrier

\subsubsection{Сводное сравнение всех методов для ETN на MoNbTaVW}

\begin{table}[H]
\centering
\small
\caption{ETN на MoNbTaVW: сводное сравнение оптимизаторов.}
\label{tab:etn-monbtavw-summary}
\begingroup
\setlength{\tabcolsep}{2.5pt}
\begin{tabular}{@{}lccc@{}}
\toprule
Метод & Validation EPA RMSE, eV/atom & Validation forces RMSE, eV/\AA & Train time, s \\
\midrule
Native BFGS & \(0.476876 \pm 2.33759\cdot 10^{-10}\) & \(0.523471 \pm 8.20466\cdot 10^{-11}\) & \(145.110 \pm 41.9128\) \\
SciPy BFGS & \(0.499952 \pm 4.27592\cdot 10^{-10}\) & \(0.529053 \pm 1.30138\cdot 10^{-9}\) & \(240.921 \pm 63.9803\) \\
SGD & \(1.72847 \pm 0.000501\) & \(0.887349 \pm 2.30430\cdot 10^{-10}\) & \(5446.38 \pm 108.408\) \\
Momentum & \(1.73013 \pm 0.000473\) & \(0.887349 \pm 2.30449\cdot 10^{-10}\) & \(5425.00 \pm 50.5405\) \\
Adam & \(0.441821 \pm 0.001756\) & \(0.678150 \pm 0.000513\) & \(5344.02 \pm 28.9838\) \\
Muon pad sqrt & \(0.428698 \pm 0.004227\) & \(0.600915 \pm 0.003125\) & \(5392.03 \pm 106.256\) \\
Muon factorization & \(0.428698 \pm 0.004227\) & \(0.600915 \pm 0.003125\) & \(5423.30 \pm 63.3786\) \\
\bottomrule
\end{tabular}
\endgroup
\end{table}

\resultfigure{fig_05_31_etn_validation_errors_monbtavw.png}{Validation errors для ETN на MoNbTaVW.}{fig:monbtavw-validation-10000}
\resultfigure{fig_05_32_etn_time_quality_monbtavw.png}{Качество и время обучения ETN на MoNbTaVW.}{fig:monbtavw-time-quality-10000}
\resultfigure{fig_05_33_etn_ensemble_spread_monbtavw.png}{Разброс ETN-результатов по ансамблю на MoNbTaVW.}{fig:monbtavw-ensemble-10000}
\resultfigure{fig_05_34_etn_loss_overview_monbtavw.png}{Обзор сглаженных loss histories для ETN на MoNbTaVW.}{fig:monbtavw-losses-overview-10000}

Сводное сравнение показывает, что по validation EPA RMSE лучшими становятся Muon pad sqrt и Muon factorization. Adam также превосходит native BFGS и SciPy BFGS по energy error, но уступает Muon. При этом по validation forces RMSE и времени обучения BFGS-варианты остаются сильнее mini-batch методов. Поэтому для MoNbTaVW итог зависит от выбранной метрики: Muon лучше по energy error, а BFGS-варианты лучше по force error и вычислительной стоимости.

\FloatBarrier

\subsection{Эксперимент с возмущением параметров ETN}

Дополнительно проверялась гипотеза о том, что периодическое добавление малых возмущений к параметрам ETN может уменьшить разброс по ансамблю или помочь выходить из неблагоприятных локальных минимумов:
\[
\theta_{t+1}
=
\theta_t
-
\eta_t g_t
+
\varepsilon_t,
\qquad
\varepsilon_t
\sim
\mathcal N(0,\sigma^2 I).
\]
Найденный набор результатов ограничен: строгий baseline CSV для сравнения отсутствует, а ETN в этой постановке практически сходился в один минимум. Поэтому этот эксперимент рассматривается как вспомогательная проверка и не используется как основное доказательство преимущества какого-либо метода.

\FloatBarrier

\section{Обсуждение результатов}
\label{sec:discussion}

\subsection{Роль квазиньютоновских методов}

BFGS и SciPy BFGS работают с полным градиентом и используют приближение обратного гессиана, то есть частично учитывают локальную кривизну функции потерь. Для обучения потенциалов это особенно важно: функционал одновременно содержит энергии и силы, а поверхность loss может быть плохо обусловленной. На MTP это даёт наиболее заметное преимущество BFGS-группы над mini-batch методами. В рассматриваемой конфигурации на AgPd минимальную validation EPA RMSE получил native BFGS, поэтому для MTP основной вывод достаточно устойчив: full-batch квазиньютоновские методы остаются сильнее first-order mini-batch оптимизаторов.

Для ETN native BFGS также сохраняет роль сильного baseline. При этом результат SciPy BFGS уже не повторяет поведение MTP: внешний BFGS-интерфейс не всегда улучшает качество и может давать большее время обучения при близкой или худшей validation error. Поэтому выводы о BFGS-вариантах нельзя переносить между MTP и ETN без отдельной проверки.

\subsection{Поведение SGD, Momentum и Adam}

SGD и Momentum в проведённых экспериментах сильно зависят от learning rate, batch size и clipping. Шум mini-batch градиента затрудняет движение по поверхности loss, а слишком крупный шаг может приводить к плато или нестабильности. Momentum сглаживает направление, но не отменяет необходимости аккуратно подбирать масштаб шага.

Adam в большинстве mini-batch сравнений оказывается сильнее SGD и Momentum. Это согласуется с его устройством: метод использует сглаженный первый момент и адаптивную нормировку по второму моменту. Для MTP на AgPd Adam даёт минимальную validation EPA RMSE среди mini-batch методов. Для ETN на AgPd он также лидирует по validation EPA RMSE, но по validation forces RMSE меньшую ошибку дают варианты Muon.

\subsection{Различие MTP и ETN}

MTP опирается на аналитически заданный набор инвариантных базисных функций. Такая структура хорошо сочетается с full-batch квазиньютоновскими методами и делает масштаб инициализации особенно заметным для SciPy BFGS. Для MTP native BFGS и SciPy BFGS близки по energy error, но native BFGS лучше по validation forces RMSE. ETN устроен иначе: он использует тензорно-сетевую параметризацию эквивариантных признаков. Это повышает гибкость модели, но меняет ландшафт оптимизации и относительное поведение методов.

Практически это означает, что оптимизатор нельзя выбирать отдельно от архитектуры потенциала. Это видно не только на SciPy BFGS, который хорошо работает для MTP при подходящей инициализации, но не становится универсально лучшим методом для ETN. Ещё более показательно ведёт себя Muon. Для MTP на AgPd варианты Muon уступают Adam и BFGS-вариантам по validation EPA RMSE. Для ETN картина меняется: на AgPd Muon почти сравнивается с Adam по energy error и превосходит его по force error, а на MoNbTaVW по validation EPA RMSE превосходит Adam и BFGS-варианты.

Такое различие хорошо согласуется со структурой моделей. Muon нормализует матричные направления обновления, поэтому для ETN с тензорно-сетевой параметризацией и большим числом матричных блоков такой тип шага выглядит более естественным, чем для MTP, где параметры связаны с явно заданными базисными инвариантами. Поэтому вывод о лучшем mini-batch методе нельзя сводить только к Adam. Adam остаётся сильным baseline, но для ETN варианты Muon становятся конкурентными, а на MoNbTaVW дают меньший средний validation EPA RMSE, чем все остальные рассмотренные методы. Даже с учётом разброса ансамбль Muon на MoNbTaVW расположен в области меньших energy errors, чем ансамбль Adam.

\subsection{Перенос выводов на MoNbTaVW}

MoNbTaVW нужен как проверка выводов на другом, более многокомпонентном датасете. В ETN-экспериментах native BFGS и SciPy BFGS остаются сильными full-batch baseline, особенно по validation forces RMSE и времени обучения. Однако после \(10000\) mini-batch шагов разница между full-batch и mini-batch подходами по energy error не только уменьшается: варианты Muon дают меньшую validation EPA RMSE, чем native BFGS и SciPy BFGS, хотя всё ещё уступают им по force error и времени обучения.

Итог по mini-batch методам становится датасет-зависимым и зависит от целевой метрики. На AgPd Adam лидирует по validation EPA RMSE, а Muon даёт меньшую force error для ETN. На MoNbTaVW варианты Muon pad sqrt и Muon factorization существенно лучше Adam, SGD и Momentum по validation EPA RMSE и становятся лучшими по energy error среди всех рассмотренных методов. Поэтому Muon нельзя рассматривать как второстепенный метод: для ETN он показывает выраженный потенциал, особенно на MoNbTaVW. При этом по force error и времени обучения BFGS-варианты пока остаются сильнее.


\section{Заключение}
\label{sec:conclusion}

В ходе работы был собран единый протокол сравнения оптимизаторов для MTP и ETN на датасетах AgPd и MoNbTaVW. В него вошли train/validation RMSE по энергиям на атом и силам, время обучения, анализ ансамблей запусков и диагностика траекторий loss.

Для проверки внешней full-batch оптимизации был реализован интерфейс \texttt{MlipWrapper}, который передаёт значение loss и градиент из MLIP-4 в \texttt{scipy.optimize.minimize}. Этот результат показывает, что Python/SciPy-оптимизаторы можно использовать для обучения MLIP-потенциалов, если доступны параметры потенциала, значение функции потерь и её градиент.

Для MTP лучший validation EPA RMSE получил native BFGS: \(0.002104 \pm 0.000096\) eV/atom. SciPy BFGS оказался близок по energy error, но хуже по forces RMSE: \(0.002333 \pm 0.000137\) eV/atom и \(0.055401 \pm 0.000759\) eV/\AA. Среди mini-batch методов на MTP лучшим по validation EPA RMSE оказался Adam: \(0.029321 \pm 0.007448\) eV/atom, однако он существенно уступил BFGS-вариантам.

Для ETN native BFGS остался сильным baseline. На AgPd native BFGS дал validation EPA RMSE 0.0138745. Среди mini-batch методов Adam дал минимальную validation EPA RMSE \(0.022675 \pm 0.000710\) eV/atom, а Muon factorization дал минимальную validation forces RMSE \(0.850381 \pm 0.005353\) eV/\AA. Поэтому для ETN на AgPd нельзя выделять единственный лучший mini-batch метод без указания метрики: Adam сильнее по energy error, а Muon --- по force error.

На MoNbTaVW вывод по energy error стал более тонким: native BFGS дал validation EPA RMSE 0.476876, SciPy BFGS -- 0.499952, Adam -- \(0.441821 \pm 0.001756\), а Muon pad sqrt и Muon factorization -- \(0.428698 \pm 0.004227\). Поэтому при 10000 mini-batch шагах варианты Muon не только существенно превосходят Adam, SGD и Momentum, но и дают меньшую validation EPA RMSE, чем BFGS-варианты. При этом по validation forces RMSE и времени обучения BFGS/SciPy остаются сильнее.

Итоговый вывод состоит в том, что оптимизатор нужно подбирать с учётом архитектуры потенциала, масштаба параметров, датасета и целевой метрики. Квазиньютоновские методы остаются основным baseline для гладкой full-batch постановки и дают лучшие ошибки по силам для ETN на MoNbTaVW. Среди mini-batch методов Adam сохраняет роль сильной опорной точки, но не является универсально лучшим вариантом. Для ETN матричный характер Muon лучше согласуется со структурой модели, чем для MTP: на AgPd Muon сравнивается с Adam и превосходит его по силам, а на MoNbTaVW превосходит Adam по validation EPA RMSE и становится лучшим методом по энергетической ошибке среди всех рассмотренных оптимизаторов. Поэтому дальнейшее исследование Muon для ETN представляется содержательным направлением.

\section{Дальнейшая работа}
\label{sec:future-work}

Дальнейшая работа естественно распадается на несколько направлений:
\begin{itemize}
    \item расширенный подбор гиперпараметров для mini-batch и квазиньютоновских методов;
    \item более осмысленное преобразование вектора параметров в матричную форму для применения Muon;
    \item проверка выводов на новых датасетах и других классах межатомных потенциалов;
    \item теоретическое исследование свойств методов оптимизации в задачах обучения машинно-обучаемых потенциалов.
\end{itemize}

\newpage
\bibliographystyle{unsrt}
\bibliography{refs}

\end{document}
